\documentclass[a4paper]{article}

% Basic packages
\usepackage[fontset=fandol]{ctex}
\usepackage{amsmath, amssymb}
\usepackage{graphicx}
\usepackage[margin=2.5cm]{geometry}
\usepackage[most]{tcolorbox}
\usepackage{etoolbox}
\usepackage{listings}
\usepackage{hyperref}
\usepackage{booktabs}
\usepackage{subcaption}
\usepackage{float}
\usepackage{tikz}
\IfFileExists{pgfplots.sty}{
    \usepackage{pgfplots}
    \pgfplotsset{compat=1.18}
}{}

% Highlight boxes
\newtcolorbox{knowledgebox}[1]{
    enhanced, colback=blue!5!white, colframe=blue!75!black, colbacktitle=blue!75!black,
    coltitle=white, fonttitle=\bfseries, title=#1, attach boxed title to top left={yshift=-2mm, xshift=2mm},
    boxrule=1pt, sharp corners
}

\newtcolorbox{importantbox}[1]{
    enhanced, colback=yellow!10!white, colframe=yellow!80!black, colbacktitle=yellow!80!black,
    coltitle=black, fonttitle=\bfseries, title=#1, sharp corners
}

\newtcolorbox{warningbox}[1]{
    enhanced, colback=red!5!white, colframe=red!75!black, colbacktitle=red!75!black,
    coltitle=white, fonttitle=\bfseries, title=#1, sharp corners
}

% Listings
\lstset{
    language=Python,
    basicstyle=\ttfamily\small,
    keywordstyle=\color{blue},
    stringstyle=\color{red!60!black},
    commentstyle=\color{green!60!black},
    breaklines=true,
    frame=single,
    numbers=left,
    numberstyle=\tiny\color{gray},
    captionpos=b,
    extendedchars=false
}

% Front-page metadata
\newcommand{\notetitle}{CS231N 2025 第4讲：神经网络与反向传播}
\newcommand{\noteauthors}{Codex Harness Worker 1}
\newcommand{\notedate}{2026-04-15}
\newcommand{\videochannel}{Stanford Online}
\newcommand{\videopublishdate}{2025-09-02}
\newcommand{\videoduration}{1:16:46}
\newcommand{\videourl}{https://www.youtube.com/watch?v=25zD5qJHYsk}
\newcommand{\videocoverpath}{cover.jpg}

\begin{document}

\begin{titlepage}
\centering
{\Large 课程笔记\par}
\vspace{1.2cm}
{\huge\bfseries \notetitle\par}
\vspace{0.8cm}
{\large \noteauthors\par}
\vspace{0.3cm}
{\large \notedate\par}
\vspace{1.2cm}

\ifdefempty{\videocoverpath}{
{\small 请在生成笔记时填入视频封面图的本地路径。\par}
}{
\includegraphics[width=0.82\textwidth,height=0.45\textheight,keepaspectratio]{\videocoverpath}\par
}

\vfill
\begin{tcolorbox}[width=0.9\textwidth, colback=black!2!white, colframe=black!60, sharp corners]
\textbf{视频作者/频道}：\videochannel\par
\textbf{发布日期}：\videopublishdate\par
\textbf{视频时长}：\videoduration\par
\textbf{视频链接}：
\ifdefempty{\videourl}{
未填写
}{
\href{\videourl}{\nolinkurl{\videourl}}
}
\end{tcolorbox}
\end{titlepage}

\tableofcontents
\newpage

%% --- 正文内容开始 --- %%


\section{从线性模型到多层感知机}
Lecture 4 的第一步，是解释为什么线性分类器不够。即使上一讲已经有了损失函数与优化器，只要模型本身仍是线性的，它就无法表达复杂的类别边界。多层感知机（multi-layer perceptron, MLP）通过在仿射变换之间插入非线性激活函数，把简单线性映射提升为高度灵活的分段线性函数族。

一个典型的两层网络可以写成
\[
 h = \sigma(W_1 x + b_1), \qquad s = W_2 h + b_2.
\]
其中，$x$ 是输入，$W_1, b_1$ 是第一层参数，$h$ 是隐藏表示，$\sigma(\cdot)$ 是非线性激活，$s$ 是输出分数。若选择 ReLU，便有 $\sigma(z) = \max(0, z)$。这里的符号解释非常重要：它说明网络先构造中间表示，再用输出层完成决策。

\begin{figure}[H]
\centering
\includegraphics[width=0.92\textwidth]{pdf_pages/page-30.png}
\caption{两层感知机通过非线性激活把线性分类器升级为分段线性函数族；“max” 或 ReLU 是表达能力跃迁的关键。}
\end{figure}

\begin{knowledgebox}{为什么非线性不可省略}
如果每一层都只是线性变换，那么多层连乘仍然等价于一层线性变换。只有引入 ReLU、sigmoid、tanh 等非线性，深度才真正增加函数表达能力。
\end{knowledgebox}

神经网络相较于线性分类器的本质升级，在于引入了可组合的非线性。
两层网络已经能表达远比线性模型更复杂的函数族。
\subsection{本章小结}

\section{深度、宽度与表示能力}
网络更深、更宽，通常意味着更大的表示容量（capacity）。但课程并没有把这个结论说成“越大越好”，而是提醒读者：更大容量同时意味着更高的过拟合风险、更难的优化问题以及更强的正则化需求。也就是说，神经网络之所以强，不是因为参数多本身，而是因为这些参数在合适的数据、目标函数和训练机制下，能够形成有用表示。

教材化理解时，可以把隐藏层看成逐步构造中间特征的过程：第一层可能对局部模式敏感，第二层组合成更抽象的部件，最后一层对类别决策负责。虽然这种解释在不同模型中并不总是字面成立，但它提供了理解深网络的第一性直觉：深度允许逐层重写表示。

深度与宽度提高了模型容量，但也带来了优化与泛化的新困难。
网络真正学到的是中间表示，而不只是最后一层分类器。
\subsection{本章小结}

\section{为什么必须高效求梯度}
一旦模型变成多层非线性结构，参数数量就迅速膨胀。若仍像上一讲那样用有限差分逐个参数近似梯度，计算代价会灾难性增长。于是课程自然转入一个核心问题：给定一个由许多简单运算组成的复杂函数，如何系统而高效地求出每个参数的梯度？

答案就是把网络看成计算图（computational graph）。前向传播负责计算每个节点的数值，反向传播则沿图结构把损失对输出的梯度逐步传回输入与参数。这个思想极其强大，因为它把求导从“手写一个巨大公式”变成了“为每个基本算子定义局部梯度规则，然后通过链式法则自动组合”。

深模型训练的关键瓶颈不是“有没有目标函数”，而是“能否高效得到所有参数的梯度”。
计算图把复杂函数拆成局部算子，为反向传播提供了统一执行框架。
\subsection{本章小结}

\section{反向传播的核心：链式法则在计算图上的执行}
课程先从标量例子入手，因为这能把链式法则讲到最透明。设某个中间变量 $q = xy$，输出为 $f = q + z$，那么损失对输入的梯度来自局部导数与上游梯度的连乘。如果已知 $\frac{\partial L}{\partial f}$，则有
\[
\frac{\partial L}{\partial x} = \frac{\partial L}{\partial f} \cdot \frac{\partial f}{\partial q} \cdot \frac{\partial q}{\partial x}, \qquad
\frac{\partial L}{\partial y} = \frac{\partial L}{\partial f} \cdot \frac{\partial f}{\partial q} \cdot \frac{\partial q}{\partial y}.
\]
其中，$\frac{\partial L}{\partial f}$ 称为上游梯度（upstream gradient），$\frac{\partial f}{\partial q}$、$\frac{\partial q}{\partial x}$ 等则是局部梯度（local gradient）。

\begin{figure}[H]
\centering
\includegraphics[width=0.92\textwidth]{pdf_pages/page-60.png}
\caption{标量计算图示例：先做前向值传播，再沿图结构逐边乘上局部梯度，这就是反向传播的工作方式。}
\end{figure}

\begin{figure}[H]
\centering
\includegraphics[width=0.92\textwidth]{pdf_pages/page-70.png}
\caption{上游梯度（upstream gradient）与局部梯度（local gradient）的乘积给出当前节点的梯度贡献。}
\end{figure}

这套机制最重要的地方在于可复用：无论整个网络多复杂，每个节点只关心两件事，第一是前向值，第二是局部梯度。前者在前向传播时缓存下来，后者在反向传播时与上游梯度相乘。把这种局部规则串起来，就得到了全局梯度。

反向传播不是神秘技巧，而是链式法则在计算图上的系统执行。
每个节点只需知道自己的前向值和局部梯度，就能参与全局求导。
\subsection{本章小结}

\section{局部梯度模式与工程实现}
课程随后总结了几个最重要的局部梯度模板。对加法节点，梯度会被原样分发到每个输入；对乘法节点，每个输入会乘上“另一个输入”的数值；对 max/ReLU 节点，只有赢得最大值的分支会接收梯度。看似琐碎的这些规则，实际上组成了几乎所有神经网络层的底层求导积木。

\begin{figure}[H]
\centering
\includegraphics[width=0.92\textwidth]{pdf_pages/page-100.png}
\caption{不同算子对应不同梯度流模式：加法分发梯度，乘法交换乘子，max 只把梯度传给获胜分支。}
\end{figure}

工程实现上，最重要的设计是 forward/backward API：前向函数返回输出并缓存必要中间量，反向函数接收上游梯度并返回对输入与参数的梯度。把每个层封装成这种接口之后，神经网络训练就变成了“缓存前向中间量 + 统一反向接口”的工程问题。

\begin{figure}[H]
\centering
\includegraphics[width=0.92\textwidth]{pdf_pages/page-110.png}
\caption{把每个层封装成 forward/backward API 后，神经网络训练就变成“缓存前向中间量 + 统一反向接口”的工程问题。}
\end{figure}

为了把“实现”说得更具体，可以把一层的训练接口写成如下伪代码：
\begin{lstlisting}[caption={单层 forward/backward 接口的伪代码}]
out, cache = layer.forward(x)
dx, dw, db = layer.backward(dout, cache)
params['W'] -= lr * dw
params['b'] -= lr * db
\end{lstlisting}

局部梯度规则与模块化实现共同构成了现代自动求导框架的基础。
理解这一节之后，再看 PyTorch 的 autograd，就不会把它当成黑箱。
\subsection{本章小结}

\section{总结与延伸}
Lecture 4 完成了从“可优化的线性模型”到“可训练的深层非线性模型”的关键跨越。它先解释了神经网络为什么需要非线性，再说明反向传播为什么能够让大规模参数学习真正可行。若没有这一讲，后续卷积网络、自注意力、多模态模型的训练都会失去基础。

从教材的组织上看，这一讲把深度学习最核心的训练机制讲清了：前向传播负责计算值，反向传播负责沿计算图传递梯度，优化器负责用梯度更新参数。这三者共同构成现代深度模型训练的主轴。下一讲进入卷积网络时，我们只是把“层的种类”换掉，而不是把训练逻辑推倒重来。

\subsection{拓展阅读}
\begin{itemize}
\item \href{http://cs231n.github.io/optimization-2}{Backprop}
\item \href{https://distill.pub/2017/momentum/}{Why Momentum Really Works}
\item \href{http://colah.github.io/posts/2015-08-Backprop/}{Backpropagation Blog}
\end{itemize}


%% --- 正文内容结束 --- %%

\end{document}
